\section{The Deployment Front-Door Rule}
\label{sec:rule}

The attribution gap calls for a limited information rule rather than a new general law of AI liability. The rule must tell an affected person whom to approach, tell enterprises what record to maintain, and tell courts when and how that record becomes available. It must also provide a prompt exit for an enterprise whose connection to the brand or technical stack did not include practical control over the asserted risk.

This Part supplies that architecture. Section A ties coverage to sectoral legislative designation. Section B makes the visible operator the first-answer gateway and defines a threshold linked to an existing legal claim. Section C gives \term{material control} a cumulative, risk-specific test. Section D separates record creation, preservation, and staged production. Section E specifies confidentiality, remedies, and safe exit. Section F consolidates the proposal into model statutory text, and Section G explains how an ex ante state-law record duty can operate alongside the Federal Rules.

\subsection{Sectoral Designation Defines the Field}
\label{sec:designation}

``High impact'' should identify a deployment category, not attach permanently to a model. A general-purpose model can support casual drafting in one setting and become part of a benefits, hiring, clinical, or physical-control process in another. The deployment supplies the protected interest, affected population, operating institution, and legally relevant control. Coverage should follow those features.

The proposed statute therefore applies only when a legislature has designated a category and specified six matters: the protected interests; the visible operator; the minimum deployment record; the retention period and material-change trigger; the claimant's threshold; and the available enforcement remedy. A legislature can make the designation directly or authorize a sectoral agency to maintain a bounded list under statutory criteria. The legislature retains the choices that determine who must answer and which interests justify compulsory recordkeeping.

Five features identify the strongest candidates. First, the deployment affects an interest already protected by substantive law, such as employment opportunity, access to public benefits, health, credit, housing, education, physical safety, or the welfare of minors. Second, the system repeatedly mediates decisions or exposure rather than producing an isolated low-consequence convenience. Third, affected people predictably observe less about the operative system than the operator and suppliers know. Fourth, the identity, version, control, and event fields can be standardized at reasonable cost. Fifth, existing doctrine can adjudicate the dispute once those facts are available. These criteria select the settings in which a record gateway unlocks an established legal regime rather than inventing a new one.

Designation also fixes the public-law interface. If a covered operator is a government institution, the legislature should specify the enforcing tribunal, available relief, applicable exhaustion route, and the treatment of sovereign or official immunity. The front door need not imply a damages action against the government. It can supply an administrative record, declaratory remedy, reviewable denial, or other access path tailored to the underlying interest. Defining that route at designation keeps the information right aligned with the sector's remedial structure.

This design avoids an unstable universal definition. Employment screening, allocation of essential public benefits, clinical diagnosis or treatment, credit, housing, education admissions, physical safety, and relational services offered to minors may justify different records and retention periods. A companion service may require age-access, relational-behavior, crisis-intervention, and continuity records. An employment tool may require input-feature, validation, selection-rate, configuration, and version records. A clinical system may require intended-use, validation-population, uncertainty, escalation, override, and follow-up records. The shared front-door architecture permits sector-specific content.

California's enacted companion-chatbot provisions demonstrate the feasibility of this form. The legislature defined a covered relational service and its operator, required specified disclosures and self-harm protocols, imposed protections for known minors and reporting obligations, supplied a private remedy, and preserved cumulative law.\lrfootnote{See \statcite{californiaSB2432025}. The enacted provisions do not contain the deployment front-door rule proposed here; they demonstrate that a legislature can designate a relational deployment, identify its operator, and attach enforceable protocol and reporting duties.} The statute did not need to declare every conversational model high impact. It regulated a deployment category and the enterprise that made it available.

A designation should also identify exclusions. Research or internal evaluation that does not produce a legally or practically consequential effect on an affected person ordinarily falls outside. So does general cloud hosting without risk-specific configuration or intervention authority. Small entities may receive simplified record forms or regulator assistance where the sector permits, but the information essential to identify a consequential decision should not disappear with enterprise size. A covered operator using open weights still controls a deployment; an upstream author who has no continuing authority may fall outside the post-release control lane.

\subsection{The Visible Operator Answers First}
\label{sec:gateway}

The rule resolves the service problem by selecting the entity an affected person can ordinarily identify. The \term{visible operator} is the enterprise or public institution that made the covered deployment available to the affected person or incorporated it into the process that produced the adverse event. It may be a companion-service provider, employer, hospital, lender, benefits agency, school, or operator of a branded application. Its gateway role follows from visibility and provenance, not a presumption that it controlled every risk.

Stage One uses an occurrence threshold rather than a miniature merits pleading. A verified notice identifies: (1) an adverse event affecting an interest protected by the law specified in the sectoral designation; (2) the covered deployment or visible operator through which the event occurred; (3) facts reasonably connecting the requested identity and event fields to that occurrence; (4) the relevant time, transaction, account, decision, or interaction with as much specificity as reasonably available; (5) the closed fields sought; and (6) a good-faith intended use of the answer to evaluate, contest, appeal, or seek relief for the event. A declaration under penalty of perjury supplies the requester's personal knowledge and marks matters stated on information reasonably available; counsel certifies proper purpose and reasonable inquiry.

That threshold demands an event and a bounded connection, but never a fact held exclusively in the deployment record. An applicant need not know whether a named feature was enabled, whether it generated a score, whether the score reached the employer, which version ran, or which supplier retained the event log. Those are first-answer fields. ``An AI harmed me'' remains insufficient because it identifies neither a covered occurrence nor a closed record. But the operator cannot defeat Stage One by demanding proof that an undisclosed feature participated in the event.

The notice invokes a statutory access right that exists outside merits discovery. It can be served before a complaint, with a complaint, or after filing. Delivery follows the electronic or registered-agent channel named in the designation; it invokes the operator's first-answer duty but is not service of process for a later action. The operator receives thirty days, subject to one short extension for a documented technical, privacy, or security need, to answer, object by field, or state that no covered deployment processed the identified event. Silence, a blanket proprietary designation, and referral to an unnamed vendor are not responses.

If the operator refuses or materially underanswers, the affected person may seek narrow enforcement against that known entity in a lawful forum. The statutory violation is complete upon coverage, a compliant request, delivery, and nonperformance; enforcement does not depend on pleading the merits fact that the first answer exists to reveal. The administering agency or a court with independent subject-matter and personal jurisdiction can order the defined answer, and a valid arbitration agreement may select the private enforcement forum. The requester bears the burden on coverage, delivery, and the occurrence threshold. The operator bears the burden on a field-specific objection based on facts within its control, including nonoperation, burden, privilege, privacy, security, or documented unavailability. Jurisdiction, immunity, preemption, and a pure legal bar remain available before access expands beyond the fixed first answer.

The gateway then supplies a \term{first answer} with two components. The \term{event receipt} identifies the application and model versions active for the event; any output, score, decision, notice, warning, transmission, use, or intervention recorded for that event; and whether a requested function did not run. The \term{custodian map} identifies entities that supplied, configured, integrated, or operated material components; the holder and legal controller of each listed record; and the roles with authority to update, restrict, roll back, suspend, or stop a relevant function. The operator preserves records within its control, identifies the retention schedule and known deletion, and activates each contractual retrieval and preservation path required before deployment.

This answer solves two practical problems. The affected person receives identity and provenance without reconstructing vendor architecture. Potential control holders receive a common event identifier and can preserve the corresponding slice. A visible operator cannot answer that procurement did not retain the version or that the supplier relationship is confidential; local provenance and a retrieval right are conditions of covered deployment. Notice to a supplier, however, creates neither statutory coverage, choice of law, personal jurisdiction, nor service of process. Direct statutory duties attach only where the enterprise and deployment fall within the Act's valid territorial and choice-of-law reach; the operator remains responsible for the local minimum answer and contractual retrieval architecture.

Stage Two begins only after that record exists. The affected person must identify a claim-specific control lane, show from Stage One that the function formed part of the event path or governed a safeguard, and show that additional nonprivileged information bears non-de-minimally on an element, defense, or remedy and is proportionate to the defined dispute. The visible operator remains free to show that its gateway role did not include the relevant lane. Visibility opens Stage One; the test in the next Section determines who remains for Stage Two.

\subsection{Material Control Requires Risk, Authority, and Information Materiality}
\label{sec:material-control}

``Ability to affect the system'' is too broad. A cloud provider, investor, data vendor, researcher, model author, and downstream institution may each have some causal connection to a deployment. The front door should reach an enterprise only when three requirements are satisfied at the relevant time.

First, \term{risk specificity}: the asserted control must connect to the risk alleged in the verified notice. General influence over corporate strategy, compute availability, or model research does not establish control over age access, hiring rank, clinical escalation, a particular tool permission, or another specified hazard.

Second, \term{practical authority}: the enterprise must have possessed an operational or legally enforceable ability to configure, constrain, evaluate, monitor, update, suspend, roll back, or terminate the relevant function. A theoretical technical capacity, an expired contract right, or influence that required another actor's voluntary cooperation is evidence but not practical authority by itself. Contractual audit, update, safety-override, and suspension rights matter because they may be usable controls; a contractual label declaring the parties independent does not decide the question.

Third, \term{information materiality}: the Stage One record must support that the function formed part of the event path or governed a safeguard, and the requested information must bear non-de-minimally on an element, defense, or remedy. An incidental contribution or a remote ability affecting every customer in the same undifferentiated way falls outside. The inquiry asks whether the requested information can materially inform adjudication. It does not require a showing that exercising the authority would have prevented the injury, changed the ultimate decision, or satisfied causation.

The requirements are cumulative. Nonexclusive factors include who selected the deployment's purpose and population; chose the model and version; configured instructions, memory, retrieval, data fields, tools, and credentials; set engagement, ranking, decision, or escalation thresholds; received evaluations, complaints, and incident signals; could issue an update or impose a restriction; retained contractual audit or safety-override rights; and could suspend, roll back, or terminate the function. The inquiry is temporal. Authority obtained after the event does not create earlier control, although later records may reveal what authority existed.

Several negative rules keep the test bounded. Investment, ordinary cloud capacity, publication of research, past employment, trademark association, and model provenance alone do not establish material control. Publishing open weights may end the original provider's ability to monitor or stop an independent deployment. A downstream modifier that removes safeguards, adds credentials, or selects a new high-impact purpose can occupy the relevant lane instead. Conversely, a provider that retains mandatory updates, remote suspension, risk monitoring, or binding deployment conditions may retain a lane despite contract language assigning ``all responsibility'' downstream.

Material control determines information obligations, not substantive liability. An enterprise can possess a control relevant to feasibility yet owe no duty to the affected person under governing law. It can receive an alert that does not constitute legally sufficient notice. The test identifies the entities and records needed to ask the merits questions accurately. The statute's \term{merits firewall} preserves that separation: front-door participation and record production create no inference on an element of the underlying claim.

\subsubsection{Adjudicating Material Control}

The test becomes useful doctrine when each requirement has an evidentiary question and an allocated burden. After the first answer, the affected person makes a prima facie showing by identifying a concrete lane connected to the asserted risk, showing that the Stage One record supports that the function formed part of the event path or governed a safeguard, and identifying requested information that bears non-de-minimally on a claim-specific element, defense, or remedy. The showing may point to an employer's enabled ranking feature, a provider's retained update authority, an application's self-harm threshold, or an integrator's workflow trigger. It concerns allocation and record relevance, not breach or causal sufficiency. An enterprise seeking termination then identifies Stage One material showing that the function did not run, no output was generated or transmitted, the asserted authority was unavailable at the relevant time, or the requested information bears no non-de-minimis relation to an element, defense, or remedy. The tribunal decides termination on that closed record or permits a bounded Stage Two inquiry when a genuine factual conflict remains.

The record duty supplies one limited burden consequence. If a covered enterprise's failure to create or retain a required Stage One field produces an ambiguity about its lane, the tribunal resolves that ambiguity against the record-duty holder only for access and termination. The consequence opens bounded production; it creates no inference of duty, breach, causation, discriminatory operation, or damages. That allocation prevents record noncompliance from becoming a safe exit while preserving the merits firewall.

\begin{table}[htbp]
\centering
\caption{Material Control as an Adjudicable Test}
\label{tab:material-control-proof}
\small
\begin{tabularx}{\textwidth}{>{\raggedright\arraybackslash}p{0.18\textwidth} >{\raggedright\arraybackslash}p{0.35\textwidth} X}
\toprule
Requirement & Evidence that supports the lane & Evidence that supports termination of the lane \\
\midrule
Risk specificity & Function map; risk ownership; feature requirements; evaluation scope; incident category; warnings; escalation assignment. & General corporate influence; unrelated research; infrastructure service without risk-specific configuration; records addressing a different feature or hazard. \\
Practical authority & Contract right; administrative permission; release or update role; audit access; monitoring feed; restriction, rollback, suspension, or termination mechanism; consistent exercise in practice. & Expired or legally unavailable right; access requiring another actor's uncommitted consent; technical ability prohibited and unused; independent modification beyond visibility or reach. \\
Information materiality & Stage One evidence that the function formed part of the event path or governed a safeguard; requested information bearing non-de-minimally on an element, defense, or remedy. & Feature inactive for the event; output never generated or transmitted; lane acquired after the event; requested information unrelated to the asserted path or claim-specific issue. \\
\bottomrule
\end{tabularx}
\end{table}

Reserved contractual rights deserve particular care. A right to obtain validation on demand, require a safety update, audit use, or suspend a feature is practical authority when the right was legally available and operationally usable. The holder need not have exercised it in the claimant's event; otherwise an enterprise could erase its information obligation by leaving a reserved safeguard idle. A generic right to terminate an entire commercial agreement may be too remote when it provides no timely, risk-specific intervention. The contract supplies evidence, while administrative permissions, actual practice, response time, technical dependencies, and past interventions show whether the right could function.

Actual power can matter when a contract is silent or inaccurate. A vendor that routinely pushes mandatory updates, disables a feature across tenants, monitors risk signals, or supplies the only version map may possess practical authority even if the agreement calls the customer solely responsible. A customer that can enable, ignore, or combine a score may control the decision lane even when the vendor hosts the data. A systems integrator with one-time access during installation may fall outside after handoff unless it retained monitoring, credentials, or change authority. The inquiry reconstructs the operative deployment rather than enforcing labels selected for another purpose.

Information materiality imposes event-path discipline without creating a causation hearing. Stage One can answer many cases: the feature was disabled; no score was generated; the output never reached the decisionmaker; the provider had no post-release connection; or the supposed controller joined the stack after the event. A lane remains when the Stage One record supports that its function formed part of that path or governed a safeguard and the requested information bears non-de-minimally on a claim-specific element, defense, or remedy. Where the event path and authority conflict, Stage Two can examine configuration history, tests, monitoring, or a representative protocol limited to that lane. No different-act, prevention, or causation showing is required.

The same test produces different allocations across architectures. In an integrated companion service, the visible operator may control access, relationship design, memory, engagement objectives, warnings, monitoring, and account restriction; an infrastructure provider offering generic compute exits. In an employment stack, the employer may control job criteria, feature activation, decision weight, review, and final disposition, while a vendor controls model design, defaults, validation, updates, and cross-customer monitoring. In an open-weights deployment, the original publisher may have no continuing control after release while a modifier and operating institution acquire the risk-specific tools, data, credentials, and stop authority. The rule follows each lane without requiring every participant to remain.

A material-control ruling terminates only the special front-door role. An entity never named on an underlying claim returns to ordinary nonparty status after completing any first-answer duty. A party already facing an independently sufficient cause of action remains subject to that claim even if it lacks the control lane alleged in the notice. This distinction gives safe exit its intended economy without using an information proceeding to decide a merits question it was not designed to reach.

\subsection{Creation, Preservation, and Production Are Separate Duties}
\label{sec:three-duties}

The front door depends on three legal functions that begin at different times.

\subsubsection{Record Creation Before Operation}

The \term{deployment-record duty} arises before a covered deployment affects the public. The visible operator has a nondelegable local duty to identify the system, version, purpose, population, material components, vendors, data sources, tools, credentials, intended limits, risk owners, release approval, event identifier, and intervention authority. During operation, it records material changes, evaluations, complaints, incidents, escalations, overrides, restrictions, rollbacks, and termination decisions. A covered control holder subject to the Act maintains the subset corresponding to its lane and provides current provenance and a response contact. If a supplier is outside the Act's direct reach, the operator still owes the local minimum record and must perform through its predeployment contract.

A divided deployment requires a retrieval architecture as well as a list of vendors. Before use, the visible operator must obtain contractual rights to retrieve the identity, version, configuration, event, and control-allocation fields assigned to an upstream holder on demand. The agreement identifies a responsive agent, stable deployment and event identifiers, preservation and response duties, lawful routes for producing customer-controlled data, successor obligations after acquisition, and the treatment of subcontractors. The operator need not warehouse a supplier's source code or unrelated customer data. It keeps a closed local record sufficient to identify what ran and a legally usable path to obtain the defined supplier fields.

That requirement builds a usable retrieval right into deployment design. Where Rule 34 control turns on a legal right to obtain documents on demand, a covered operator can create that right before deployment instead of litigating it after harm. Where privacy law or a customer agreement places legal control with the operating institution, the contract can authorize a coordinated response rather than allowing the host and customer to direct the requester to each other. Operator notice alone does not create personal jurisdiction over a foreign or out-of-state supplier.\lrfootnote{See \casecite[264--68]{bristolMyers2017}; \casecite[359--71]{fordMotor2021}. The model statute therefore distinguishes a control holder directly subject to the Act from a supplier reached through the operator's contract and local record.} A domestic operator preserves the minimum record locally and conditions covered use on a reachable agent and retrieval duty.

The designation supplies the detail. A statute need not require indefinite retention of every prompt or every user's data. It should define event categories, sampling where appropriate, privacy minimization, and a retention period calibrated to the protected interest and ordinary limitations period. High-impact decisions affecting employment or benefits may require the inputs, output, version, governing policy, notice, and review event for each affected person while permitting aggregation or deletion of unrelated interaction content. Relational deployments involving minors may require crisis detections and interventions while minimizing intimate content unrelated to the asserted event.

\subsubsection{Preservation After a Defined Trigger}

Retention and litigation preservation are related but distinct. The statute fixes a baseline period before any dispute. A verified front-door notice, regulator demand, filed claim, or other event making a dispute reasonably foreseeable suspends ordinary deletion for the risk-specific record. The operator must issue a hold to relevant custodians and notify identified control holders. The notice should not freeze every model and customer record. Its scope follows the event, versions, period, user or decision, control lanes, and asserted risk.

This statutory trigger creates a stable bridge to existing preservation law. It also protects defendants: a dated notice and documented hold define what had to be preserved and when. The record can show that a missing item was outside scope, deleted before any duty attached, or available from another source.

\subsubsection{Staged Production After the Threshold}

Production proceeds in stages. Stage One is the first answer described above: identity and active versions; event-specific run, output, notice, and intervention fields; suppliers, custodians, and legal controllers; listed operational and stopping authority; and preservation, deletion, and retrieval status. Its purpose is to identify the deployment and test the control lanes.

Stage Two follows only when the affected person uses Stage One to make a prima facie showing of a claim-specific control lane, information materiality, and proportional need. It can include relevant evaluations and validation; complaints and incident trends; configuration history; release and change approvals; rejected or conditional mitigations; escalation and response; and records needed to test a proposed alternative or defense. Time, version, affected population, feature, and risk category bound the request.

Source code, model weights, unrestricted training data, security-sensitive controls, and unrelated users' content occupy a final tier. A court should order them only when the requesting party shows that a material issue cannot reasonably be tested through less intrusive evidence such as versioned documentation, validation results, representative testing, counterfactual interrogation, interrogatories, an expert protocol, or in camera review.\lrfootnote{See \bluecite[1076--79]{stewartBeyondExplanation2026}. Counterfactual interrogation tests the behavior of a known decision-time version; Stage One supplies the identity and preservation predicates needed to direct that tool.} The distinction protects legitimate secrecy without converting a trade secret into a barrier to the legal process that governs its consequential use.

\subsection{Protection, Remedies, and Safe Exit}
\label{sec:remedies}

Confidentiality is part of the rule's design. Production remains limited to nonprivileged matter. Courts may use redaction, pseudonymization, data minimization, attorneys'-eyes-only access, secure expert environments, in camera review, and protective orders for trade secrets, research, security, medical information, and personal data.\lrfootnote{See \statcite{frcp26cTradeSecret}.} The operator must still identify that a protected record exists, its custodian, date, version, and asserted basis for withholding. A privilege log separates a lawful protection from an unreviewable claim that the record is proprietary.

The statute should specify graduated consequences.

\begin{enumerate}[leftmargin=2.2em]
\item \term{Late or incomplete identification}: an order enforcing the statutory first answer and reasonable expenses caused by noncompliance. For a claim created or governed by the same enactment, the limitations period is suspended from service of a compliant request until a fixed period after compliance or final denial. Other state and federal claims retain their own limitations and tolling law.
\item \term{Failure to create or retain the minimum record}: the administrative penalty, corrective record order, compliance monitoring, or statutory-access remedy specified by the sectoral enactment. The consequence addresses the independent deployment obligation; it does not become a prescribed inference on the underlying injury.
\item \term{Loss after a preservation trigger}: an order to identify the loss and available substitute sources, followed by measures under the governing forum's preservation law. For electronically stored information that should have been preserved in anticipation or conduct of federal litigation and cannot be restored or replaced, Rule 37(e) supplies the federal standard. The state statute should not prescribe a competing litigation sanction for the same loss.
\item \term{Intentional concealment or falsification}: specified civil penalties, enforcement expenses, and remedies for the independent access violation. Litigation misconduct remains subject to the tribunal's governing authority.
\item \term{The merits firewall}: compliance, noncompliance, and front-door participation create no inference on the merits questions enumerated in Section 8. Governing law determines whether a record omission has any separate relevance to an independently pleaded claim.
\end{enumerate}

Safe exit is the reciprocal limit. After Stage One, an enterprise may move to terminate its front-door role by showing that it lacked material control over the asserted risk during the relevant period. The claimant may use the first-answer record to contest that showing. An entity not otherwise sued returns to ordinary nonparty status. A party remains only when an independently sufficient claim, a valid subpoena, or regulatory authority supplies a basis for continued participation. The gateway obligation therefore ends at the point where the control record shows that another lane or no AI-specific lane is relevant.

\subsection{Model Statutory Text}
\label{sec:model-text}

The following text consolidates the rule as a default model. A sectoral enactment would identify the covered field, retention clock, administering authority, and ordinary civil-penalty ceiling while preserving the sequence and merits firewall.

\begin{quote}
\small
\textbf{Section 1. Designated Deployment and Covered Enterprises.} This Act applies only to an AI deployment category designated by the enacting legislature and identified by covered function, protected interests, visible operator, affected persons, minimum record, retention period, material-change threshold, administering authority, enforcement route, and civil-penalty ceiling. A \emph{covered control holder} is an enterprise to which this State's law validly applies that operates, provides, or continuously supports the identified covered deployment, or that has contractually assumed the response duties defined here. Before designation, the legislature or authorized agency shall publish a federal-overlap determination identifying applicable federal actors, statutes, regulations, exclusive enforcement provisions, express or conflict-preemption issues, excluded actors and fields, the relation between remedies, and the claims to which state tolling may lawfully apply.

\textbf{Section 2. Deployment Record and Retrieval Architecture.} Before making a covered deployment available or incorporating it into an affected process, the visible operator shall create and maintain a local record identifying the system and versions; purpose and affected population; material model, data, retrieval, tool, and credential components; deployment and event identifiers; suppliers and integrators; risk and release owners; evaluation and approval; material limitations and unresolved risk decisions; incident and escalation channels; and authority to update, restrict, roll back, suspend, or terminate. Each covered control holder shall maintain the corresponding portion and provide current provenance and a response contact. The operator shall obtain a contractual right to retrieve the defined supplier fields, require preservation and response, and carry those duties through subcontracting, acquisition, and service termination. These operator duties are nondelegable. The designation supplies the retention period and material-change rule.

\textbf{Section 3. Stage One Notice.} An affected person may deliver to the visible operator's designated agent a sworn declaration or declaration under penalty of perjury identifying: (a) an adverse event affecting an interest protected by the law named in the designation; (b) the covered deployment or visible operator through which the event occurred; (c) facts reasonably connecting the requested identity and event fields to that occurrence; (d) the event, transaction, account, application, decision, or period with as much specificity as reasonably available; (e) the closed fields sought; and (f) a good-faith intended use of the answer to evaluate, contest, appeal, or seek relief for the event. Stage One does not require a fact held exclusively in the deployment record. Counsel, if any, shall certify proper purpose and reasonable inquiry. Delivery invokes the first-answer duty but is not service of process for an enforcement action.

\textbf{Section 4. First Answer and Preservation.} Within thirty days after delivery, the operator shall provide an event receipt and custodian map. The receipt identifies the application and model versions active for the event; any requested output, score, decision, notice, warning, transmission, use, or intervention; and whether a named function did not run. The map identifies each entity that supplied, configured, integrated, or operated a material component; the custodian and legal controller of each field; and each entity with update, restriction, rollback, suspension, or termination authority. The operator shall preserve its risk-specific record, identify known deletion or unavailability, and activate its contractual retrieval and preservation rights. Each covered control holder independently subject to this Act shall preserve and acknowledge its assigned lane. For any other supplier, the operator shall satisfy the retrieval obligation through Section 2's contract and local record. Operator notice alone creates neither statutory coverage, choice of law, service, personal jurisdiction, nor consent. One extension of fifteen days may be obtained for a stated technical, privacy, or security need.

\textbf{Section 5. Material Control.} An enterprise remains for Stage Two only upon a prima facie showing that, at the relevant time: (a) its function was specifically connected to the asserted risk; (b) it possessed practical authority through configuration, constraint, evaluation, monitoring, update, restriction, rollback, suspension, termination, or a legally enforceable equivalent; and (c) the Stage One record supports that the function formed part of the event path or governed a safeguard, and the requested information bears non-de-minimally on an element, defense, or remedy. This Section requires no showing that exercising the authority would have prevented the injury, changed the decision, or satisfied causation. Investment, general cloud service, publication, research, brand association, provenance, or theoretical technical capacity alone is insufficient.

\textbf{Section 6. Objection, Burdens, and Enforcement.} The operator may object by field and shall state facts supporting noncoverage, nonoperation, burden, privilege, privacy, security, or documented unavailability. A proprietary label or referral to an unnamed supplier is insufficient. A timely, specific objection stays disclosure of only the contested field pending prompt neutral review; every uncontested field remains due. An affected person may seek expedited enforcement in a court of competent jurisdiction or another lawful forum. The requester bears the burden on coverage, delivery, and the Stage One occurrence threshold. The operator bears the burden on a field-specific objection based on facts within its control. After the first answer, the requester identifies a prima facie control lane; an enterprise seeking termination identifies the Stage One record defeating risk specificity, practical authority, or information materiality. An ambiguity caused by that enterprise's violation of Section 2 is resolved against it only for access and termination and creates no merits inference.

\textbf{Section 7. Stage Two and Protection.} Further nonprivileged production may be obtained by agreement or ordered in a pending proceeding only when Stage One supports a claim-specific control lane, including that the function formed part of the event path or governed a safeguard; the requested information bears non-de-minimally on an element, defense, or remedy; and production is necessary and proportionate. Production shall be limited by time, version, affected population, feature, and asserted risk. Source code, weights, sensitive data, raw content, expressive deliberation, and security controls require written findings that less intrusive evidence cannot reasonably resolve a material issue. The owner of contested protected material and any materially affected third party shall receive notice and an opportunity for prompt neutral review before disclosure. The tribunal shall protect privilege, privacy, trade secrets, research, association, expression, and security through minimization, redaction, pseudonymization, limited access, secure expert review, in camera review, or another lawful safeguard. In a pending federal action, the Federal Rules govern Stage Two production and protection.

\textbf{Section 8. Merits Firewall and Termination.} Compliance, noncompliance, notice, and participation create no presumption of duty, breach, defect, discriminatory operation, agency, causation, damages, or joint liability. Governing law supplies every element and defense. An enterprise that completes Section 4 and defeats a Section 5 requirement terminates its front-door role. An entity not otherwise subject to an independently sufficient claim resumes nonparty status. A valid subpoena, regulatory authority, or another source of law may require continued participation on its own terms.

\textbf{Section 9. Remedies and Limitations.} The lawful forum may issue an expedited compliance order and award reasonable actual expenses caused by an incomplete, late, or evasive answer. The administering authority may order creation or correction of the minimum record, compliance monitoring, and a civil penalty calibrated to duration, affected persons, record scope, and culpability up to the ceiling fixed in the sectoral designation. Knowing falsification, concealment, or repeated violation permits recovery of enforcement expenses and an enhanced penalty up to twice that ceiling. Bare nonperformance creates no private statutory damages; this Act creates no damages for the underlying adverse event and no merits presumption. Measures for lost electronically stored information in federal litigation remain governed by Federal Rule of Civil Procedure 37(e). Only a claim created or expressly governed by this Act is tolled from delivery of a compliant request until sixty days after compliance or final denial.

\textbf{Section 10. Nonwaiver, Forum, Standing, and Federal Procedure.} The duties to create, retain, and provide the minimum record are nonwaivable as to affected persons, arise from covered deployment, and exist whether a dispute is litigated, arbitrated, or never filed. A valid arbitration agreement may select the private enforcement forum under the Federal Arbitration Act; the first-answer clock continues independently and is not a condition precedent to commencing or continuing arbitration. An agreement to which the administering authority is not a party does not restrict public enforcement brought in that authority's own name. Nothing in this Act creates Article III standing, personal jurisdiction, or subject-matter jurisdiction; overrides immunity or privilege; binds a nonsignatory to arbitration; or alters federal rules governing pleading, service, joinder, discovery, amendment, protective orders, injunctions, and sanctions. Courts and arbitrators administer the duties created here; they do not derive them from ordinary discovery authority.

\textbf{Section 11. Territorial Reach, Federal Supremacy, and Severability.} An upstream enterprise has a direct duty only when it purposefully supplied, configured, operated, or continuously supported the identified deployment for covered use in this State, or another lawful basis for legislative authority exists, and application of this State's law satisfies territorial and choice-of-law rules. Coercive relief additionally requires valid service and an independent basis for personal jurisdiction. Appointment of the operator's response agent authorizes delivery under Section 3 only and creates no consent to general jurisdiction. This Act imposes no requirement to the extent federal law preempts it, requires no disclosure prohibited by federal privacy, security, communications, or confidentiality law, creates no claim solely to enforce a federal duty reserved to the United States, and does not regulate the United States absent congressional authorization and an applicable waiver. A tribunal shall sever a preempted actor, field, remedy, or application when the remaining designation can operate consistently with the legislature's stated purpose.
\end{quote}

The model text makes the statutory architecture legible. The regulated acts are operating a covered deployment without the minimum record and refusing the resulting access right, not unsafe operation in the abstract. The trigger is tied to a legislatively specified legal interest, an adverse event, and closed deployment facts. Control is risk specific. Access is staged. Confidential information receives protection. A record holder can terminate its special role. The merits remain where ordinary law placed them.

\subsection{Why the First Answer Is a Substantive Access Right Rather Than Accelerated Discovery}
\label{sec:erie}

A state front-door statute regulates covered operation rather than rewriting federal litigation. The minimum-record duty attaches when the operator makes a designated deployment available or incorporates it into an affected process. The first-answer duty attaches to a qualifying request whether a lawsuit is anticipated, pending, arbitrated, or never filed. Its content is fixed by the designation---identity, event, provenance, custody, and control fields---rather than by the relevance scope of a particular complaint. Refusal completes an independent statutory violation. Claim-specific discovery begins only at Stage Two.

That design matters because state law cannot evade a Federal Rule by relabeling a pleading or discovery requirement. Federal courts have rejected state presuit certificate requirements when they conflict with the Federal Rules' answer to what a federal complaint must contain.\lrfootnote{See, e.g., \casecite[516--24]{pledger2021}.} The front door asks a different question. It creates a record during operation and gives the affected person a transaction-linked access entitlement before suit. A federal complaint enforcing that entitlement remains governed by Rule 8; any additional production remains governed by Rules 26, 34, and 45.

The right and its enforcement forum therefore remain separate. A state statute does not manufacture federal subject-matter jurisdiction. An affected person ordinarily proceeds in the designated state forum, invokes diversity if its requirements are met, asserts a federal right if Congress enacts one, or seeks supplemental jurisdiction when the state access claim forms part of a case already within federal jurisdiction.\lrfootnote{See \statcite{federalSubjectMatterJurisdiction}.} Once the dispute is in federal court, federal rules govern pleading, service, discovery, amendment, protection, and sanctions.

Article III adds a separate limit to a federal access action. A legislature's grant of a cause of action informs the inquiry but does not make every statutory violation concrete. In \emph{TransUnion}, the plaintiffs before the Court could not establish concrete informational injury without alleging downstream adverse effects; \emph{Spokeo} likewise separates a legal violation from concrete injury in fact.\lrfootnote{See \casecite[425--31, 440--42]{transUnion2021}; \casecite[339--43]{spokeo2016}.} Informational-standing decisions found concreteness where plaintiffs were personally denied information that law required to be disclosed and the information bore directly on a protected use.\lrfootnote{See \casecite[373--74]{havens1982}; \casecite[449--50]{publicCitizen1989}; \casecite[21--25]{akins1998}.} A federal front-door plaintiff should therefore plead that denial impaired the ability to identify, contest, appeal, or seek relief for the specified adverse event, or caused another concrete effect such as reasonable actual expense. A bare delay untethered to such use may fail in federal court even though the legislature created a cause of action. State courts remain governed by their own constitutions and standing law; Article III does not require a state forum to close whenever a parallel federal action would lack standing.\lrfootnote{See \casecite[617--18]{asarco1989}.}

The Rules Enabling Act framework then makes direct collision the first Erie question. \emph{Hanna} applies a valid Federal Rule when it answers the question before the court. \emph{Burlington Northern} rejected a state appellate penalty that occupied the same field as a federal rule governing the consequence at issue. \emph{Walker}, by contrast, applied a state service-linked limitations rule because Rule 3 did not answer when the state limitations period was tolled.\lrfootnote{See \casecite[464--65, 468, 471--74]{hanna1965}; \casecite[4--8]{burlingtonNorthern1987}; \casecite[749--52]{walker1980}; \statcite{rulesEnablingAct}.} The fixed deployment and access duties sit on the \emph{Walker} side only to the extent they regulate operation and create an entitlement that exists outside litigation. They neither advance Rule 26's timetable nor enlarge Rule 34 control or Rule 45 subpoena power.

The allocation described in Section~\ref{sec:ordinary-procedure} follows by function. State law can require a version, provenance, control, and event record years before a dispute; create a nonwaivable request; authorize an action for refusal; and suspend limitations for a claim the enactment creates or governs. Rule 8 controls the federal enforcement complaint. Rules 26, 34, and 45 control Stage Two in federal court. Rule 26(c) supplies protective orders. Rule 15 controls amendment and relation back. The state statute gives the parties an existing object and access entitlement; federal procedure supplies the litigation machinery.

Preservation requires the same separation. The deployment statute defines baseline retention and can penalize noncreation or premature deletion as an operating violation. A qualifying request can create an independent statutory hold. Rule 37(e), however, governs federal litigation measures when electronically stored information that should have been preserved in anticipation or conduct of the action is lost and cannot be restored or replaced. The Rule supplies prejudice and intent standards; it does not create the operating duty.\lrfootnote{See \statcite{frcp37e}.} The model text therefore provides regulatory and access remedies for the deployment violation while leaving federal adverse-inference and case-management measures to Rule 37(e).

Limitations follows legislative authority. A state can suspend the period for a state claim it creates or governs while the operator supplies the first answer. It cannot promise tolling for an independent federal cause of action or automatic relation back against a later-named party. Congress can pair a federal front door with federal tolling where national uniformity is warranted. In either form, the access period gives the affected person time to use the answer while preserving Rule 15's role in a later federal case.

The resulting division is concrete. Deployment law keeps a consequential system identifiable and grants a defined access right. Federal procedure governs the federal enforcement action and any later merits claim. The same allocation preserves the merits disagreement identified in the Introduction: the access right supplies the deployment and record on which the governing liability rule can operate, while the merits firewall neither selects that rule nor alters its elements. The first answer arrives with the missing record; it does not arrive as a substitute code of civil procedure.
